\documentclass[11pt]{article}

\usepackage[margin=1in]{geometry}
\usepackage{amsmath, amssymb}
\usepackage{algorithm}
\usepackage{algpseudocode}
\usepackage{hyperref}

\title{Upper Confidence Bound (UCB)}
\author{Reinforcement Learning Notes}
\date{}

\begin{document}
\maketitle

\section{The Multi-Armed Bandit Problem}

We have $d$ options (``arms''), e.g. $d$ ad versions to show to users. At each
round $n = 1, \dots, N$ we must choose one arm $i \in \{1, \dots, d\}$ to
play. Playing arm $i$ returns a reward $r_i(n) \in \{0, 1\}$ (e.g.\ whether
the user clicked the ad).

Each arm $i$ has an unknown, fixed expected reward $\mu_i$ (its true
click-through rate). The goal is to maximize cumulative reward over $N$
rounds, which requires balancing:

\begin{itemize}
  \item \textbf{Exploitation} --- keep choosing the arm that looks best so far.
  \item \textbf{Exploration} --- try other arms to check whether they are
        actually better, since estimates from few trials are unreliable.
\end{itemize}

UCB resolves this exploration/exploitation trade-off deterministically, by
choosing at each round the arm with the highest \emph{upper confidence
bound} on its expected reward.

\section{Notation}

At the start of round $n$, for each arm $i$ define:

\begin{itemize}
  \item $N_i(n)$: number of times arm $i$ has been selected in the rounds
        before $n$.
  \item $R_i(n)$: sum of rewards obtained from arm $i$ so far.
  \item $\bar{r}_i(n) = \dfrac{R_i(n)}{N_i(n)}$: empirical average reward
        of arm $i$ (exploitation term).
\end{itemize}

\section{Confidence Interval and Upper Bound}

By Hoeffding's inequality, the true mean $\mu_i$ lies within
\[
  \bar{r}_i(n) - \Delta_i(n) \;\le\; \mu_i \;\le\; \bar{r}_i(n) + \Delta_i(n)
\]
with high probability, where the confidence radius is
\[
  \Delta_i(n) = \sqrt{\frac{3}{2} \cdot \frac{\log(n)}{N_i(n)}}.
\]

Note that $\Delta_i(n)$:
\begin{itemize}
  \item \textbf{shrinks} as $N_i(n)$ grows --- the more we've tried arm $i$,
        the more confident we are in $\bar{r}_i(n)$ (exploitation dominates).
  \item \textbf{grows slowly} with $n$ (the round number) --- arms we
        haven't tried in a while regain some uncertainty, encouraging
        exploration.
\end{itemize}

The \emph{upper confidence bound} for arm $i$ at round $n$ is
\[
  \text{UCB}_i(n) = \bar{r}_i(n) + \Delta_i(n).
\]

This is an optimistic estimate of $\mu_i$: the best-case value consistent
with the data observed so far. This is the ``optimism in the face of
uncertainty'' principle.

\section{Algorithm}

\begin{algorithm}[h]
\caption{Upper Confidence Bound}
\begin{algorithmic}[1]
\State $N_i \gets 0$, $R_i \gets 0$ for all $i = 1, \dots, d$
\For{$n = 1, \dots, N$}
  \For{$i = 1, \dots, d$}
    \If{$N_i > 0$}
      \State $\bar{r}_i \gets R_i / N_i$
      \State $\Delta_i \gets \sqrt{\tfrac{3}{2}\log(n)/N_i}$
      \State $\text{UCB}_i \gets \bar{r}_i + \Delta_i$
    \Else
      \State $\text{UCB}_i \gets +\infty$ \Comment{force every arm to be tried once}
    \EndIf
  \EndFor
  \State play arm $i^\ast = \arg\max_i \text{UCB}_i$
  \State observe reward $r$; update $N_{i^\ast} \gets N_{i^\ast} + 1$,
         $R_{i^\ast} \gets R_{i^\ast} + r$
\EndFor
\end{algorithmic}
\end{algorithm}

The $+\infty$ initialization (\texttt{1e400} in the notebook code) guarantees
every arm is sampled at least once during the first $d$ rounds, since an
untried arm always beats any finite UCB value.

\section{Correspondence with the Implementation}

The notebook (\texttt{upper\_confidence\_bound.ipynb}) implements exactly
this loop over the \texttt{Ads\_CTR\_Optimisation.csv} dataset, with
$N = 10000$ rounds and $d = 10$ ads:

\begin{center}
\begin{tabular}{ll}
\textbf{Math} & \textbf{Code} \\
\hline
$N_i(n)$ & \texttt{numbers\_of\_selections[i]} \\
$R_i(n)$ & \texttt{sums\_of\_rewards[i]} \\
$\bar{r}_i(n)$ & \texttt{average\_reward} \\
$\Delta_i(n)$ & \texttt{delta\_i} \\
$\text{UCB}_i(n)$ & \texttt{upper\_bound} \\
$\arg\max_i \text{UCB}_i(n)$ & \texttt{ad} (tracked via \texttt{max\_upper\_bound}) \\
\end{tabular}
\end{center}

The final histogram of \texttt{ads\_selected} shows that, as $n$ grows, the
algorithm converges to selecting the ad with the highest true CTR far more
often than the others --- the exploration term shrinks for the best arm
as its estimate becomes reliable, while inferior arms are selected less
and less.

\section{Key Takeaways}

\begin{itemize}
  \item UCB is a \emph{deterministic} strategy: no randomness is used to
        pick an arm, unlike Thompson Sampling.
  \item It requires updating estimates after every single round (cannot be
        batched/deferred), since $N_i(n)$ and $n$ both change per round.
  \item It has strong theoretical regret bounds: cumulative regret grows
        only logarithmically in $N$.
\end{itemize}

\end{document}
